\section{Validación del modelo}

Cualquier modelo predictivo necesita ser validado para observar como es su rendimiento en diferentes conjuntos de datos  y determinar si la precisión del modelo es contante todas fuentes de datos similares o no.


Esto nos ayuda a detectar algún \texttt{problema del exceso de ajuste (over-fitting)}, en el que el modelo se ajusta
muy bien en un conjunto de datos, pero no encaja bien en otro conjunto de datos. Un método común es crear un modelo diviendo los datos en categorías de \texttt{entrenamiento/prueba}. Otro método es
\texttt{validación cruzada de $k$ iteraciones}, sobre la cual aprenderemos más en el capítulo posterior.

\subsection{División entre datos de entrenamiento y prueba}

Idealmente, este paso debería hacerse justo al inicio del proceso de modelado para que
no haya sesgos de muestreo en el modelo; en otras palabras, el modelo debería funcionar
bien incluso para un conjunto de datos que tiene las mismas variables de predicción, pero sus medias y
las varianzas son muy diferentes sobre de las que se ha construido el modelo.


Esto puede
suceder porque el conjunto de datos en el que se basa el modelo (entrenamiento o capacitación) y el de
que se aplica (prueba) puede provenir de diferentes fuentes. Una forma más robusta de
hacer esto es un proceso llamado la validación cruzada de $k$ iteraciones, sobre la cual hablaremos en
detalle en un momento.


Veamos cómo podemos dividir el conjunto de datos disponible en el conjunto de datos de entrenamiento y prueba
y aplicar el modelo al conjunto de datos de prueba para obtener otros resultados:

[,]{\texttt{crossValidation.py}}
\begin{lstlisting}[language=Python]
	import pandas as pd
	import statsmodels.formula.api as smf
	
	advert = pd.read_csv("./dataBases/Advertising.csv")
	model3=smf.ols(formula='Sales~TV+Radio',data=advert).fit()
	sales_pred=model3.predict(advert[['TV','Radio']])
	print(sales_pred.head())
	
	print(model3.summary())
	
	import numpy as np
	
	N = len(advert)
	arr =  np.arange(N)
	np.random.shuffle(arr)
	check = arr < 0.8*N
	training=advert[check].copy()
	testing=advert[~check].copy()
\end{lstlisting}


Vamos a crear un modelo para entrenar los datos y problemaar el rendimiento del modelo en
datos de prueba. Creemos el único modelo que funciona mejor (lo hemos encontrado ya),
el que tiene variables de TV y radio, como variables de predicción:

[,]{}
\begin{lstlisting}[language=Python]
	import statsmodels.formula.api as smf
	model5=smf.ols(formula='Sales~TV+Radio',data=training).fit()
	print(model5.summary())
\end{lstlisting}

[,]{\texttt{print(model3.summary())}}

\begin{lstlisting}[language=Python]
	OLS Regression Results
	==============================================================================
	Dep. Variable:                  Sales   R-squared:                       0.897
	Model:                            OLS   Adj. R-squared:                  0.896
	Method:                 Least Squares   F-statistic:                     859.6
	Date:                Sun, 12 Nov 2017   problema (F-statistic):           4.83e-98
	Time:                        22:40:26   Log-Likelihood:                -386.20
	No. Observations:                 200   AIC:                             778.4
	Df Residuals:                     197   BIC:                             788.3
	Df Model:                           2
	Covariance Type:            nonrobust
	==============================================================================
	coef    std err          t      P>|t|      [0.025      0.975]
	------------------------------------------------------------------------------
	Intercept      2.9211      0.294      9.919      0.000       2.340       3.502
	TV             0.0458      0.001     32.909      0.000       0.043       0.048
	Radio          0.1880      0.008     23.382      0.000       0.172       0.204
	==============================================================================
	Omnibus:                       60.022   Durbin-Watson:                   2.081
	problema(Omnibus):                  0.000   Jarque-Bera (JB):              148.679
	Skew:                          -1.323   problema(JB):                     5.19e-33
	Kurtosis:                       6.292   Cond. No.                         425.
	==============================================================================
	
	Warnings:
	[1] Standard Errors assume that the covariance matrix of the errors is correctly specified.
\end{lstlisting}

[,]{}

\begin{lstlisting}[language=Python]
	OLS Regression Results
	==============================================================================
	Dep. Variable:                  Sales   R-squared:                       0.906
	Model:                            OLS   Adj. R-squared:                  0.904
	Method:                 Least Squares   F-statistic:                     753.7
	Date:                Sun, 12 Nov 2017   problema (F-statistic):           3.22e-81
	Time:                        22:40:26   Log-Likelihood:                -299.60
	No. Observations:                 160   AIC:                             605.2
	Df Residuals:                     157   BIC:                             614.4
	Df Model:                           2
	Covariance Type:            nonrobust
	==============================================================================
	coef    std err          t      P>|t|      [0.025      0.975]
	------------------------------------------------------------------------------
	Intercept      2.8708      0.318      9.035      0.000       2.243       3.498
	TV             0.0448      0.001     30.797      0.000       0.042       0.048
	Radio          0.1959      0.009     22.803      0.000       0.179       0.213
	==============================================================================
	Omnibus:                       11.888   Durbin-Watson:                   2.158
	problema(Omnibus):                  0.003   Jarque-Bera (JB):               13.175
	Skew:                          -0.696   problema(JB):                      0.00138
	Kurtosis:                       2.807   Cond. No.                         438.
	==============================================================================
	
	Warnings:
	[1] Standard Errors assume that the covariance matrix of the errors is correctly specified.
\end{lstlisting}


La mayoría de los parámetros del modelo, como la intercepción, las estimaciones de coeficientes y $R^2$ son
muy similares.


La diferencia en el estadístico$-F$ se puede atribuir a un conjunto de datos más pequeño. Cuanto menor sea el conjunto de datos, mayor será el valor de \texttt{SSD} y menor será el valor de
el término $(n-p-1)$ en la fórmula del estadístico$-F$; ambos contribuyen a la disminución en el valor estadístico$-F.$


El modelo puede reescribirse como sigue:
\begin{align}
	\texttt{Ventas} \sim 2.86 + 0.04*\texttt{TV} + 0.17*\texttt{Radio}
\end{align}

[,]{}
Ahora, predigamos los valores de las ventas para los valores de prueba:

\begin{lstlisting}[language=Python]
	sales_pred=model5.predict(training[['TV','Radio']])
	sales_pred
\end{lstlisting}


El valor \texttt{RSE} para esta predicción en el conjunto de datos de prueba puede ser calculadas usando el siguiente pedazo de código:

[,]{}
\begin{lstlisting}[language=Python]
	testing['sales_pred']=2.86 + 0.04*testing['TV'] + 0.17*testing['Radio']
	n = len(testing)
	p = 2
	testing['SSD']=(testing['Sales']-testing['sales_pred'])**2
	SSD=testing.sum()['SSD']
	RSE=np.sqrt(SSD/(n-p-1))
	salesmean=np.mean(testing['Sales'])
	error=RSE/salesmean
	print(RSE,salesmean,error)
	##2.33080393856 14.527500000000003 0.160440814907
\end{lstlisting}


El valor \texttt{RSE} resulta ser $2.54$ sobre una venta promedio (en el conjunto de datos) de $14.80$, lo cual es un error del $17\%$.


Podemos ver que el modelo no se generaliza muy bien en el conjunto de datos de prueba, ya que el valor \texttt{RSE} para el mismo modelo es diferente en los dos casos.


Implica un cierto grado de exceso de ajuste cuando tratamos de construir el modelo basado en todo el conjunto de datos.


El valor \texttt{RSE} con
la división de pruebas de entrenamiento, aunque un poco mayor, es más confiable y replicable.

\subsection{Validación cruzada de $k$ iteraciones ($k$-fold cross-validation)}

La \emph{validación cruzada de $k$ iteraciones} es una técnica más robusta que la simple división entrenamiento/prueba. En lugar de hacer una sola partición, dividimos los datos en $k$ subconjuntos (folds) de tamaño aproximadamente igual, y realizamos $k$ iteraciones donde en cada una:

\begin{enumerate}
	\item Usamos $k-1$ folds para entrenar el modelo.
	\item Usamos el fold restante para probar el modelo.
	\item Calculamos las métricas de rendimiento (ej: $R^2$, $RSE$).
\end{enumerate}

Al final, promediamos las métricas obtenidas en las $k$ iteraciones para obtener una estimación más estable y confiable del rendimiento del modelo.

\begin{observacion}
	Ventajas de la validación cruzada $k$-fold:
	\begin{itemize}
		\item \textbf{Menor varianza:} Al promediar sobre $k$ particiones, reducimos la variabilidad debida a una partición específica.
		\item \textbf{Uso eficiente de datos:} Cada observación se usa exactamente una vez para prueba y $k-1$ veces para entrenamiento.
		\item \textbf{Detección de sobreajuste:} Si el rendimiento en entrenamiento es consistentemente mejor que en prueba, hay evidencia de sobreajuste.
	\end{itemize}
\end{observacion}

\subsubsection{Implementación en Python}

\begin{lstlisting}[language=Python]
import pandas as pd
import numpy as np
from sklearn.model_selection import KFold
from sklearn.linear_model import LinearRegression
from sklearn.metrics import mean_squared_error, r2_score

# Cargar datos
advert = pd.read_csv("./dataBases/Advertising.csv")
X = advert[['TV', 'Radio']]
Y = advert['Sales']

# Configurar validación cruzada
k = 5  # número de folds
kf = KFold(n_splits=k, shuffle=True, random_state=42)

# Almacenar métricas
r2_scores = []
rse_scores = []

# Realizar validación cruzada
for fold, (train_idx, test_idx) in enumerate(kf.split(X), 1):
    # Dividir datos
    X_train, X_test = X.iloc[train_idx], X.iloc[test_idx]
    Y_train, Y_test = Y.iloc[train_idx], Y.iloc[test_idx]
    
    # Entrenar modelo
    model = LinearRegression()
    model.fit(X_train, Y_train)
    
    # Predecir en conjunto de prueba
    Y_pred = model.predict(X_test)
    
    # Calcular métricas
    r2 = r2_score(Y_test, Y_pred)
    mse = mean_squared_error(Y_test, Y_pred)
    rse = np.sqrt(mse)
    
    r2_scores.append(r2)
    rse_scores.append(rse)
    
    print(f'Fold {fold}: R² = {r2:.4f}, RSE = {rse:.4f}')

# Resultados agregados
print(f'\n=== Resultados de Validación Cruzada {k}-Fold ===')
print(f'R² promedio: {np.mean(r2_scores):.4f} ± {np.std(r2_scores):.4f}')
print(f'RSE promedio: {np.mean(rse_scores):.4f} ± {np.std(rse_scores):.4f}')
\end{lstlisting}

\subsubsection{Interpretación de resultados}

Los resultados de la validación cruzada nos proporcionan:

\begin{itemize}
	\item \textbf{Métrica promedio:} Estimación del rendimiento esperado del modelo en datos nuevos.
	\item \textbf{Desviación estándar:} Medida de la estabilidad del modelo. Una desviación estándar pequeña indica que el modelo es robusto a diferentes particiones de datos.
	\item \textbf{Rendimiento por fold:} Permite identificar si hay folds problemáticos (ej: con outliers o datos atípicos).
\end{itemize}

\begin{observacion}
	\textbf{Elección de $k$:}
	\begin{itemize}
		\item $k = 5$ o $k = 10$ son valores comúnmente usados y proporcionan un buen balance entre sesgo y varianza.
		\item $k$ pequeño (ej: $k = 2$ o $k = 3$): Mayor sesgo, menor varianza, más rápido.
		\item $k$ grande (ej: $k = n$, llamado \emph{leave-one-out}): Menor sesgo, mayor varianza, más lento.
	\end{itemize}
\end{observacion}

\subsubsection{Comparación de modelos con validación cruzada}

La validación cruzada es especialmente útil para comparar diferentes modelos:

\begin{lstlisting}[language=Python]
from sklearn.model_selection import cross_val_score

# Modelo 1: Solo TV
X1 = advert[['TV']]
r2_model1 = cross_val_score(LinearRegression(), X1, Y, cv=5, scoring='r2')

# Modelo 2: TV + Radio
X2 = advert[['TV', 'Radio']]
r2_model2 = cross_val_score(LinearRegression(), X2, Y, cv=5, scoring='r2')

# Modelo 3: TV + Radio + Newspaper
X3 = advert[['TV', 'Radio', 'Newspaper']]
r2_model3 = cross_val_score(LinearRegression(), X3, Y, cv=5, scoring='r2')

print('Comparación de modelos (R² promedio ± desviación estándar):')
print(f'Modelo 1 (TV): {np.mean(r2_model1):.4f} ± {np.std(r2_model1):.4f}')
print(f'Modelo 2 (TV+Radio): {np.mean(r2_model2):.4f} ± {np.std(r2_model2):.4f}')
print(f'Modelo 3 (TV+Radio+Newspaper): {np.mean(r2_model3):.4f} ± {np.std(r2_model3):.4f}')
\end{lstlisting}

\subsubsection{Validación cruzada estratificada}

Cuando trabajamos con datos que tienen características importantes que deben preservarse en cada fold (ej: proporción de clases en clasificación, rangos de valores en regresión), podemos usar \emph{validación cruzada estratificada}:

\begin{lstlisting}[language=Python]
from sklearn.model_selection import StratifiedKFold

# Crear categorías basadas en cuartiles de Y
Y_cuartiles = pd.qcut(Y, q=4, labels=False)

# Validación cruzada estratificada
skf = StratifiedKFold(n_splits=5, shuffle=True, random_state=42)

r2_stratified = []
for train_idx, test_idx in skf.split(X, Y_cuartiles):
    X_train, X_test = X.iloc[train_idx], X.iloc[test_idx]
    Y_train, Y_test = Y.iloc[train_idx], Y.iloc[test_idx]
    
    model = LinearRegression()
    model.fit(X_train, Y_train)
    Y_pred = model.predict(X_test)
    
    r2_stratified.append(r2_score(Y_test, Y_pred))

print(f'R² estratificado: {np.mean(r2_stratified):.4f} ± {np.std(r2_stratified):.4f}')
\end{lstlisting}

\subsection{Resumen de técnicas de validación}

\begin{center}
	\begin{tabular}{|l|l|l|l|}\hline
		\textbf{Método} & \textbf{Ventajas} & \textbf{Desventajas} & \textbf{Uso recomendado} \\\hline
		Train/Test split & Rápido, simple & Alta varianza & Exploración inicial \\\hline
		$k$-fold CV & Robusto, eficiente & Más lento & Evaluación final \\\hline
		LOOCV ($k=n$) & Casi insesgado & Muy lento, alta varianza & Muestras pequeñas \\\hline
		CV estratificada & Preserva distribución & Más compleja & Datos desbalanceados \\\hline
	\end{tabular}
\end{center}

\begin{observacion}
	\textbf{Recomendación práctica:} Use validación cruzada de 5 o 10 folds como método estándar para evaluar modelos de regresión. Reserve la simple división entrenamiento/prueba solo para exploración inicial o cuando tenga conjuntos de datos muy grandes.
\end{observacion}
